% How the BERT of question 3 is fine-tuned: data, labels, tags, settings, runs.
% The numbers come from results/training/*.json, written by experiments/train_bert.py (2026-11-24).
%   cd docs/reports/train && latexmk -pdf training_procedure.tex
\documentclass[10pt,a4paper]{article}
\usepackage[T1]{fontenc}
\usepackage{lmodern}
\usepackage[margin=2cm]{geometry}
\usepackage{microtype}
\usepackage{amsmath}
\usepackage{booktabs}
\usepackage{tabularx}
\usepackage{enumitem}
\usepackage{caption}
\usepackage[hidelinks]{hyperref}
\captionsetup{font=small,labelfont=bf,skip=4pt}
\setlist{itemsep=1pt,topsep=3pt}
\setlength{\parindent}{0pt}
\setlength{\parskip}{4pt}
\newcommand{\file}[1]{\texttt{#1}}
\newcolumntype{L}{>{\raggedright\arraybackslash}X}

\begin{document}

\begin{center}
{\Large\bfseries PIIGuard: how BERT is fine-tuned}\\[3pt]
{\small Training procedure of question 3 \textperiodcentered{} CSED490E, PIIGuard team
\textperiodcentered{} November 24, 2026}
\end{center}

Question 3 asks whether training on noisy text wins back what noise costs a detector, and what it
costs on clean text (proposal, Appendix B). This document gives the training procedure in full.
The data it uses is described in \file{docs/reports/data/data\_generation.pdf}; the results are
in the Q2/Q3 report (\file{docs/reports/q2\_q3/}).

\section{The two versions}

\begin{tabularx}{\linewidth}{@{}llL@{}}
\toprule
Version & Model folder & Training data \\
\midrule
(a) clean & \file{models/bert-clean-s<seed>} & \file{ai4privacy\_train\_clean}: 115,533 records \\
(b) clean + noisy & \file{models/bert-noisy-s<seed>} & \file{ai4privacy\_train\_clean} + \file{ai4privacy\_train\_noisy}: 231,066 records (each record and one mixed-noise copy of it) \\
\bottomrule
\end{tabularx}

Everything else is the same: the starting model, the settings, the number of steps and the seeds
(0, 1, 2). The only difference is the noisy copies, so a difference between (a) and (b) comes
from them. Six models in all. The models are not in git (about 110 MB each); each run's record is in git
(\file{results/training/<model>.json}).

\section{Starting model}

\begin{itemize}
\item \file{gravitee-io/bert-small-pii-detection}, revision \file{f8c27a85c51c0168f07b9dcf00265bf0a4097939}:
  the BERT of the benchmark (28.5M parameters), already trained by Gravitee to tag 25 PII classes
  named as in Presidio.
\item Its classification head is kept as it is: the same 51 tags (\file{O}, and \file{B-} and
  \file{I-} for each of the 25 classes) and the same tag ids. So
  \file{baselines/bert\_detector.py} reads a fine-tuned model exactly like the original, with the
  same class mapping; only the weights differ.
\item Its tokenizer is uncased WordPiece: every text is lowercased before the model sees it.
\end{itemize}

\section{From gold spans to token tags}

Our labels are taught as the model's own classes:

\begin{center}
\begin{tabular}{@{}ll@{\hspace{3em}}ll@{}}
\toprule
Our label & BERT class & Our label & BERT class \\
\midrule
NAME & PERSON & ADDRESS & LOCATION \\
EMAIL & EMAIL\_ADDRESS & DATE (every date) & DATE\_TIME \\
PHONE & PHONE\_NUMBER & CARD\_NUMBER & CREDIT\_CARD \\
\bottomrule
\end{tabular}
\end{center}

\begin{itemize}
\item \textbf{Dates.} The training data uses the comparison-C labels: every date is \file{DATE},
  taught as \file{DATE\_TIME}. ai4privacy has no birth-date label; leaving its dates untagged
  would teach the model that dates are not PII. At test time the detector maps \file{DATE\_TIME}
  to \file{DATE\_OF\_BIRTH} (six labels) or \file{DATE} (\file{\_dates} versions), as for the
  original model.
\item \textbf{Other tokens} get \file{O}, so the model also learns that what ai4privacy labels
  otherwise (cities, zip codes, titles, \ldots) is not one of our labels.
\item \textbf{Tagging} (\file{bio\_labels} in \file{experiments/train\_bert.py}). The text is
  tokenized with character offsets, cut at 256 tokens. Each token takes the gold span in which its
  first character falls: \file{B-<class>} if it is the first token of that span, \file{I-<class>}
  for the next ones, \file{O} outside every span. \file{[CLS]} and \file{[SEP]} get $-100$ and are
  left out of the loss.
\item \textbf{Length.} 256 tokens is the window \file{BertDetector} reads. ai4privacy records are
  short (median 29 tokens); 1 record of 115,533 is longer and is cut.
\end{itemize}

\section{Settings}

\begin{tabularx}{\linewidth}{@{}lL@{}}
\toprule
Setting & Value \\
\midrule
Trainer & Hugging Face \file{Trainer} (transformers 5.16.1, torch 2.14.1) \\
Optimizer & AdamW (the \file{Trainer}'s default), weight decay 0.01 \\
Learning rate & $5\times10^{-5}$, linear warm-up over the first 6\% of the steps (650), then linear decay to 0 \\
Batch & 32 records per step, padded per batch (\file{DataCollatorForTokenClassification}) \\
Steps & $\lceil 3 \times 115{,}533 / 32 \rceil = 10{,}832$ for \textbf{both} versions: three passes over the clean training set. Version (b), with twice the records, sees each record 1.5 times \\
Precision & bf16 \\
Seeds & \file{seed} = \file{data\_seed} = 0, 1 or 2. No layer is newly initialized, so the seed sets only the order of the batches and the dropout \\
Evaluation during training & dev clean and dev noisy (5,000 records each), every 1,805 steps, batch 128; logged only \\
Model kept & the last step. The dev sets never choose a checkpoint: the noisy dev set would favour version (b) \\
Tuning & none. The settings were fixed before the runs and are the same for both versions \\
\bottomrule
\end{tabularx}

The settings are written once in \file{experiments/train\_bert.py} (\file{SETTINGS}); the command
line only chooses the version and the seed.

\section{Runs}

The six runs were made at the same time, one per GPU, on November 24, 2026.

\begin{center}
\begin{tabular}{@{}llrrrr@{}}
\toprule
Model & GPU & Seconds & Training loss & Dev loss, clean & Dev loss, noisy \\
\midrule
\file{bert-clean-s0} & RTX 4090 & 113 & 0.0161 & 0.0124 & 0.1314 \\
\file{bert-clean-s1} & L40S & 118 & 0.0162 & 0.0119 & 0.1330 \\
\file{bert-clean-s2} & RTX 4090 & 136 & 0.0161 & 0.0124 & 0.1361 \\
\midrule
\file{bert-noisy-s0} & RTX 4090 & 112 & 0.0258 & 0.0138 & 0.0267 \\
\file{bert-noisy-s1} & RTX 4090 & 106 & 0.0256 & 0.0126 & 0.0253 \\
\file{bert-noisy-s2} & RTX 4090 & 148 & 0.0258 & 0.0126 & 0.0254 \\
\bottomrule
\end{tabular}
\captionof{table}{Training loss: mean over the run. Dev losses: at the last step.}
\end{center}

Version (b) has a higher training loss because half of its records are noisy, which is harder.

\begin{center}
\begin{tabular}{@{}llrrrrrr@{}}
\toprule
Model & Dev set & step 1,805 & 3,610 & 5,415 & 7,220 & 9,025 & 10,830 \\
\midrule
(a) clean & clean & 0.0173 & 0.0135 & 0.0127 & 0.0128 & 0.0127 & 0.0122 \\
(a) clean & noisy & 0.0945 & 0.1048 & 0.1143 & 0.1250 & 0.1352 & 0.1335 \\
(b) clean + noisy & clean & 0.0210 & 0.0162 & 0.0146 & 0.0129 & 0.0134 & 0.0130 \\
(b) clean + noisy & noisy & 0.0411 & 0.0319 & 0.0285 & 0.0264 & 0.0269 & 0.0258 \\
\bottomrule
\end{tabular}
\captionof{table}{Dev loss during training, mean over the three seeds (\file{results/training/summary.md}).}
\end{center}

The clean-only version gets worse on the noisy dev set as training goes on (0.095 to 0.134), in
all three seeds; version (b) improves on both dev sets. Token-level loss is not the benchmark's
score: the versions are compared with strict span F2 on the test sets.

\section{Outputs and evaluation}

\begin{itemize}
\item \file{models/bert-<version>-s<seed>/} holds the weights, the configuration, the tokenizer
  and \file{training.json}. Not in git.
\item \file{results/training/bert-<version>-s<seed>.json}: the same record, in git: version, seed,
  base model and revision, settings, number of steps and records, SHA-256 of the four data files,
  every logged loss, final dev losses, seconds, date, GPU, package versions.
\item \textbf{Evaluation.} Each model is run on the clean and noisy test sets like the other
  detectors, by \file{BertDetector} with \file{-{}-model}: windows of 256 tokens overlapping by 64,
  spans grouped by the transformers pipeline, the same span clean-up and scoring. A version's
  score is the mean of its three seeds.
\end{itemize}

\section{Reproducing}

{\small
\begin{verbatim}
uv run python experiments/prepare_train.py      # data/train/ (docs/reports/data/)
for seed in 0 1 2; do
  CUDA_VISIBLE_DEVICES=$seed     uv run python experiments/train_bert.py clean --seed $seed &
  CUDA_VISIBLE_DEVICES=$((seed+3)) uv run python experiments/train_bert.py noisy --seed $seed &
done; wait                                      # about 2 minutes, six GPUs
uv run python experiments/run_benchmark.py bert --model models/bert-noisy-s0 --name bert-noisy-s0 \
    --versions clean noisy_typo noisy_space noisy_case noisy_punct noisy_number noisy_mixed
                                                # every model, and again with --dates
uv run python experiments/score.py
uv run python experiments/summarize_data.py     # results/training/summary.md
\end{verbatim}
}

\file{train\_bert.py -{}-max-steps 50} makes a short run to check the set-up. A run made again on
another GPU or package version can differ slightly in its last digits; the seed spread above shows
how much.

\section{Limits}

\begin{itemize}
\item One training domain (ai4privacy), with its own span conventions: both versions learn its
  short addresses (building number and street), which costs them on Gretel.
\item One set of settings, not tuned; three seeds per version.
\item The tokenizer is uncased, so lowercase noise cannot be learned or tested on this model.
\end{itemize}

\end{document}
